\documentclass[11pt]{article}

\usepackage[a4paper,margin=2.2cm]{geometry}
\usepackage[T2A]{fontenc}
\usepackage[utf8]{inputenc}
\usepackage[russian]{babel}
\usepackage{amsmath,amssymb}
\usepackage{hyperref}
\hypersetup{colorlinks=true,urlcolor=blue}
\sloppy

\setlength{\parindent}{1.25cm}
\setlength{\parskip}{0pt}
\linespread{1.0}

\begin{document}

\section{Полная математическая постановка задачи}

Рассматривается задача восстановления нейтронного спектра по откликам детекторов. Для каждого образца заданы:
\begin{equation}
\mathbf{x} = (x_1,\dots,x_n)^\top \in \mathbb{R}^n,
\qquad
\mathbf{y} = (y_1,\dots,y_m)^\top \in \mathbb{R}^m,
\end{equation}
где в основной постановке $n=10$ --- число каналов детекторов $Q_1,\dots,Q_{10}$, а $m=60$ --- число энергетических бинов спектра.

Требуется построить отображение
\begin{equation}
\hat{\mathbf{y}} = f_{\theta}(\mathbf{x}),
\end{equation}
где $f_{\theta}$ --- многовыходная ANFIS-модель с параметрами $\theta$.

В текущей основной версии используется двухэтапная схема обучения:
\begin{enumerate}
    \item построение базовой модели ANFIS на данных, указанных в \texttt{dataset.train\_data};
    \item дифференцируемая донастройка той же модели на реальных измерениях с SHAP-регуляризацией и тихоновской регуляризацией гладкости спектра.
\end{enumerate}

Ниже приведены именно те формулы, которые соответствуют текущей реализации в проекте \texttt{bonner-shap-reconstruction} для основной версии \texttt{Two-Stage SHAP + Tikhonov}.

\section{Нормализация по сумме сигналов}

В основном конфиге включена нормализация по сумме входных каналов (\texttt{normalize\_sum=true}). Для $i$-го образца вводится
\begin{equation}
S_i = \sum_{j=1}^{n} x_{ij}.
\end{equation}
Если $S_i = 0$, то в коде он заменяется малой положительной величиной $\varepsilon_{\mathrm{sum}} > 0$.

Нормализованные входы и выходы имеют вид
\begin{equation}
\mathbf{x}_i' = \frac{\mathbf{x}_i}{S_i},
\qquad
\mathbf{y}_i' = \frac{\mathbf{y}_i}{S_i}.
\end{equation}

Обучение выполняется в нормализованном пространстве, а при необходимости денормализация прогноза производится по формуле
\begin{equation}
\hat{\mathbf{y}}_i = S_i \hat{\mathbf{y}}_i'.
\end{equation}

Такая схема сохраняет форму спектра и убирает доминирование образцов с большой абсолютной интенсивностью.

\section{Многовыходная ANFIS-модель}

Используется Sugeno-type ANFIS с $R$ правилами. В основной конфигурации $R=50$, а тип функций принадлежности --- Gaussian.

Для правила $r \in \{1,\dots,R\}$ задаются посылки
\begin{equation}
\text{IF } x_1 \text{ is } A_{r1}
\text{ AND } \dots \text{ AND } x_n \text{ is } A_{rn},
\end{equation}
и линейное многовыходное следствие
\begin{equation}
\mathbf{z}_r(\mathbf{x}) = A_r \mathbf{x} + \mathbf{b}_r,
\qquad
A_r \in \mathbb{R}^{m \times n},
\quad
\mathbf{b}_r \in \mathbb{R}^{m}.
\end{equation}

Для Gaussian membership function можно записать
\begin{equation}
\mu_{rj}(x_j) = \exp\!\left(-\frac{(x_j-c_{rj})^2}{2\sigma_{rj}^2 + \varepsilon}\right),
\end{equation}
где $c_{rj}$ и $\sigma_{rj}$ --- центр и ширина функции принадлежности.

Сила срабатывания правила вычисляется как произведение принадлежностей:
\begin{equation}
 w_r(\mathbf{x}) = \prod_{j=1}^{n} \mu_{rj}(x_j).
\end{equation}
Нормализованный вес правила:
\begin{equation}
\bar{w}_r(\mathbf{x}) = \frac{w_r(\mathbf{x})}{\sum_{s=1}^{R} w_s(\mathbf{x}) + \varepsilon}.
\end{equation}
Тогда итоговый прогноз по всем $m$ выходам равен
\begin{equation}
\hat{\mathbf{y}} = f_{\theta}(\mathbf{x}) = \sum_{r=1}^{R} \bar{w}_r(\mathbf{x}) \, \mathbf{z}_r(\mathbf{x}).
\end{equation}

Таким образом, каждый выходной бин спектра является гладкой взвешенной комбинацией линейных локальных моделей, а нелинейность задаётся системой нечетких правил.

\section{Двухэтапное обучение}

\subsection{Этап 1: базовая ANFIS-модель}

На первом этапе обучается базовая модель без SHAP-штрафов. Если обозначить обучающий набор первого этапа через
\begin{equation}
\mathcal{D}_{\mathrm{base}} = \{(\mathbf{x}_i, \mathbf{y}_i)\}_{i=1}^{N_{\mathrm{base}}},
\end{equation}
то решается задача
\begin{equation}
\theta^{(0)} = \arg\min_{\theta} \mathcal{L}_{\mathrm{main}}(\theta; \mathcal{D}_{\mathrm{base}}).
\end{equation}

В коде этот этап реализован через \texttt{BioAnfisRegressor} с глобальной оптимизацией \texttt{OriginalPSO}; в основном конфиге используются 150 эпох PSO и популяция 100 частиц.

\subsection{Этап 2: SHAP + Tikhonov донастройка}

На втором этапе параметры инициализируются значением $\theta^{(0)}$, после чего модель дообучается на реальных данных. Пусть
\begin{equation}
\mathcal{D}_{\mathrm{real}} = \{(\mathbf{x}_i, \mathbf{y}_i)\}_{i=1}^{N_{\mathrm{real}}}
\end{equation}
--- набор реальных измерений. В текущем основном пайплайне он делится в отношении
\begin{equation}
60\% : 20\% : 20\%
\end{equation}
на подмножества
\begin{equation}
\mathcal{D}_{\mathrm{real}}^{\mathrm{train}},
\qquad
\mathcal{D}_{\mathrm{real}}^{\mathrm{val}},
\qquad
\mathcal{D}_{\mathrm{real}}^{\mathrm{test}}.
\end{equation}

Для набора из 375 образцов это соответствует разбиению $225/75/75$. Именно на финальных $75$ реальных образцах считаются основные метрики качества.

\section{Основная функция ошибки восстановления}

Для мини-батча
\begin{equation}
\mathcal{B} = \{(\mathbf{x}_i, \mathbf{y}_i)\}_{i=1}^{B}
\end{equation}
основная ошибка восстановления спектра задаётся как средняя квадратичная ошибка:
\begin{equation}
\mathcal{L}_{\mathrm{main}} = \frac{1}{Bm} \sum_{i=1}^{B} \left\| f_{\theta}(\mathbf{x}_i) - \mathbf{y}_i \right\|_2^2.
\end{equation}

Здесь множитель $1/m$ появляется потому, что в реализации используется \texttt{torch.nn.MSELoss()} с усреднением по всем элементам тензора предсказаний. Если расписать по выходным бинам, то
\begin{equation}
\mathcal{L}_{\mathrm{main}} = \frac{1}{Bm} \sum_{i=1}^{B} \sum_{k=1}^{m} (\hat{y}_{ik} - y_{ik})^2.
\end{equation}

Эта величина является главной задачей оптимизации и остаётся в итоговом функционале на обоих этапах.

\section{Градиентная proxy-оценка важности признаков}

Во второй стадии используется дифференцируемая proxy-аппроксимация глобальной важности входных признаков. Для батча предсказаний
\begin{equation}
\hat{\mathbf{Y}} = f_{\theta}(\mathbf{X}) \in \mathbb{R}^{B \times m}
\end{equation}
сначала вводится скаляризованный выход для каждого образца:
\begin{equation}
 s_i = \frac{1}{m} \sum_{k=1}^{m} \hat{y}_{ik}.
\end{equation}

Далее для каждого объекта и признака вычисляется градиент
\begin{equation}
 g_{ij} = \frac{\partial s_i}{\partial x_{ij}}.
\end{equation}

Пер-образцовая важность признака определяется как
\begin{equation}
 u_{ij} = |g_{ij}| \, |x_{ij}|.
\end{equation}

После усреднения по батчу получаем глобальную важность на уровне батча:
\begin{equation}
 v_j = \frac{1}{B} \sum_{i=1}^{B} u_{ij}.
\end{equation}

Чтобы избежать вырожденных нулевых координат, в реализации применяется отсечение снизу:
\begin{equation}
 \tilde{v}_j = \max\!\left(v_j, 10^{-6} \max_{\ell} v_{\ell}\right).
\end{equation}

Затем выполняется $L_1$-нормировка:
\begin{equation}
 p_j = \frac{\tilde{v}_j}{\sum_{\ell=1}^{n} \tilde{v}_{\ell} + \varepsilon},
 \qquad
 \mathbf{p} = (p_1,\dots,p_n)^\top.
\end{equation}

Вектор $\mathbf{p}$ и является основной дифференцируемой proxy-оценкой глобальной важности признаков, с которой работают все SHAP-компоненты.

\section{Аппроксимация ``истинной'' SHAP-важности в коде}

Согласованность в проекте вычисляется не через полный перебор коалиций и не через независимые SHAP-значения для каждого образца, а через более дешёвую аппроксимацию на уровне батча.

Для текущего батча вводится средний образец
\begin{equation}
\bar{\mathbf{x}} = \frac{1}{B} \sum_{i=1}^{B} \mathbf{x}_i.
\end{equation}

В качестве базового вектора используется средний вектор по SHAP-обучающей выборке:
\begin{equation}
\mathbf{b} = \frac{1}{N_{\mathrm{train}}} \sum_{i=1}^{N_{\mathrm{train}}} \mathbf{x}_i.
\end{equation}

Для каждого признака $j$ строится маскированный вектор
\begin{equation}
\bar{\mathbf{x}}^{(j \leftarrow b_j)} = (\bar{x}_1,\dots,\bar{x}_{j-1}, b_j, \bar{x}_{j+1},\dots,\bar{x}_n)^\top.
\end{equation}

Пусть
\begin{equation}
\bar{f}(\mathbf{x}) = \frac{1}{m} \sum_{k=1}^{m} f_{\theta,k}(\mathbf{x})
\end{equation}
--- средний по энергетическим бинам отклик модели. Тогда маскирующая важность на уровне батча определяется как
\begin{equation}
 \psi_j = \left| \bar{f}(\bar{\mathbf{x}}) - \bar{f}(\bar{\mathbf{x}}^{(j \leftarrow b_j)}) \right|.
\end{equation}

После этого выполняется нормировка:
\begin{equation}
 q_j = \frac{\max(\psi_j,0)}{\sum_{\ell=1}^{n} \max(\psi_{\ell},0) + \varepsilon},
 \qquad
 \mathbf{q} = (q_1,\dots,q_n)^\top.
\end{equation}

Если сумма $\sum_j \psi_j$ оказывается вырожденной, в коде используется запасной равномерный вектор $q_j = 1/n$.

Важно, что $\mathbf{q}$ не пересчитывается на каждом батче заново: при \texttt{use\_true\_shap=true} он кэшируется и обновляется раз в $K$ батчей, где по умолчанию $K=10$.

\section{Полная SHAP-регуляризация второй стадии}

Итоговый SHAP-штраф состоит из четырёх компонент:
\begin{equation}
\mathcal{L}_{\mathrm{SHAP}} = w_{\mathrm{cons}} R_{\mathrm{cons}} + w_{\mathrm{spar}} R_{\mathrm{spar}} + w_{\mathrm{faith}} R_{\mathrm{faith}} + w_{\mathrm{stab}} R_{\mathrm{stab}}.
\end{equation}

В основной версии веса $w_{\cdot}$ вычисляются адаптивно. Ниже подробно расписана каждая компонента.

\subsection{Компонента согласованности}

Согласованность сравнивает дифференцируемую proxy-важность $\mathbf{p}$ с маскирующей аппроксимацией $\mathbf{q}$. Используются два члена: среднеквадратическое отклонение и дивергенция Йенсена--Шеннона.

Сначала вводится
\begin{equation}
\operatorname{MSE}(\mathbf{p},\mathbf{q}) = \frac{1}{n} \sum_{j=1}^{n} (p_j - q_j)^2.
\end{equation}

Далее
\begin{equation}
\mathbf{m} = \frac{1}{2}(\mathbf{p}+\mathbf{q}),
\end{equation}
\begin{equation}
\operatorname{KL}(\mathbf{p}\|\mathbf{m}) = \sum_{j=1}^{n} p_j \log \frac{p_j}{m_j + \varepsilon},
\qquad
\operatorname{KL}(\mathbf{q}\|\mathbf{m}) = \sum_{j=1}^{n} q_j \log \frac{q_j}{m_j + \varepsilon},
\end{equation}
и
\begin{equation}
\operatorname{JS}(\mathbf{p}\|\mathbf{q}) = \frac{1}{2}\operatorname{KL}(\mathbf{p}\|\mathbf{m}) + \frac{1}{2}\operatorname{KL}(\mathbf{q}\|\mathbf{m}).
\end{equation}

Тогда согласованность в точной реализации равна
\begin{equation}
R_{\mathrm{cons}} = \frac{\operatorname{MSE}(\mathbf{p},\mathbf{q}) + \lambda_{\mathrm{JS}}\operatorname{JS}(\mathbf{p}\|\mathbf{q})}{1 + \mathcal{L}_{\mathrm{main}}},
\qquad
\lambda_{\mathrm{JS}} = 0.5.
\end{equation}

Деление на $1+\mathcal{L}_{\mathrm{main}}$ реализует precision-aware scaling: когда модель ещё неточна, давление интерпретируемости ослабляется.

\subsection{Компонента разреженности}

Главная идея --- поощрять концентрацию важности на небольшом числе детекторов, не сводя распределение к жёсткому one-hot режиму.

Нормированная энтропия важности записывается как
\begin{equation}
H(\mathbf{p}) = -\frac{1}{\log n} \sum_{j=1}^{n} p_j \log(p_j + \varepsilon).
\end{equation}

Также используется коэффициент Джини. Если $p_{(1)} \le \dots \le p_{(n)}$ --- компоненты $\mathbf{p}$ после сортировки по возрастанию, то
\begin{equation}
G(\mathbf{p}) = \frac{2 \sum_{j=1}^{n} j p_{(j)}}{n \sum_{j=1}^{n} p_{(j)} + \varepsilon} - \frac{n+1}{n}.
\end{equation}

Пусть $G_{\star}$ --- целевой уровень концентрации. В основной реализации по умолчанию используется
\begin{equation}
G_{\star} = 0.3.
\end{equation}

Штраф на недостаточную концентрацию:
\begin{equation}
R_{\mathrm{gini}} = \max(0, G_{\star} - G(\mathbf{p}))^2.
\end{equation}

Далее строится базовый штраф разреженности
\begin{equation}
R_{\mathrm{spar}}^{\mathrm{base}} = H(\mathbf{p}) + R_{\mathrm{gini}}.
\end{equation}

Итоговая формула с precision-aware scaling имеет вид
\begin{equation}
R_{\mathrm{spar}} = R_{\mathrm{spar}}^{\mathrm{base}}
\left(
\rho \frac{1}{1 + \mathcal{L}_{\mathrm{main}}} + (1-\rho)
\right),
\qquad
\rho = 0.7.
\end{equation}

То есть при хорошем качестве модели штраф на энтропию и недостающий Джини действует сильнее, а при плохом --- мягче.

\subsection{Компонента достоверности объяснений}

В данной реализации faithfulness считается не по полному многомерному вектору выхода, а по его скаляризации через среднее по энергетическим бинам. Это важно, потому что именно такая формула используется в коде.

Для каждого объекта батча определим
\begin{equation}
 s_i = \frac{1}{m} \sum_{k=1}^{m} \hat{y}_{ik},
 \qquad
 s_i^{(0)} = \frac{1}{m} \sum_{k=1}^{m} f_{\theta,k}(\mathbf{0}).
\end{equation}

Здесь базовая точка выбрана нулевым вектором
\begin{equation}
\mathbf{x}^{(0)} = \mathbf{0} \in \mathbb{R}^n.
\end{equation}

Градиент уже был определён как $g_{ij} = \partial s_i / \partial x_{ij}$. Тогда первая вариационная аппроксимация изменения отклика имеет вид
\begin{equation}
 a_i = \sum_{j=1}^{n} g_{ij}(x_{ij} - x_j^{(0)}) = \sum_{j=1}^{n} g_{ij}x_{ij}.
\end{equation}

Штраф достоверности объяснений вычисляется по формуле
\begin{equation}
R_{\mathrm{faith}} = \frac{1}{1 + \mathcal{L}_{\mathrm{main}}}
\cdot
\frac{1}{B} \sum_{i=1}^{B} \left( (s_i - s_i^{(0)}) - a_i \right)^2.
\end{equation}

Следовательно, этот член минимизирует расхождение между реальным изменением среднего отклика модели и его линейной градиентной аппроксимацией относительно нулевой базовой точки.

\subsection{Компонента стабильности}

В текущем тренере стабильность задаётся не через добавление шума к входам внутри функционала потерь, а через дисперсию важностей внутри батча.

Напомним, что
\begin{equation}
 u_{ij} = |g_{ij}|\,|x_{ij}|.
\end{equation}

Средняя карта важности по батчу:
\begin{equation}
 \bar{u}_j = \frac{1}{B} \sum_{i=1}^{B} u_{ij}.
\end{equation}

Тогда стабильность определяется как
\begin{equation}
R_{\mathrm{stab}} = \frac{1}{1 + \mathcal{L}_{\mathrm{main}}}
\cdot
\frac{1}{Bn} \sum_{i=1}^{B} \sum_{j=1}^{n} (u_{ij} - \bar{u}_j)^2.
\end{equation}

Таким образом, модель поощряется к тому, чтобы образцы внутри батча имели согласованную структуру объяснений по каналам детекторов.

\subsection{Адаптивные веса компонент SHAP}

Если включён режим \texttt{use\_adaptive\_weights=true} (а в основной конфигурации именно так), веса четырёх компонент пересчитываются динамически на каждом батче.

Пусть
\begin{equation}
\boldsymbol{\ell} = (R_{\mathrm{cons}}, R_{\mathrm{spar}}, R_{\mathrm{faith}}, R_{\mathrm{stab}}).
\end{equation}

Сначала строятся обратные веса по величине потерь:
\begin{equation}
\tilde{w}_j = \frac{1}{\ell_j + \varepsilon},
\qquad
w_j^{(0)} = \frac{\tilde{w}_j}{\sum_{r} \tilde{w}_r}.
\end{equation}

Далее, если основная ошибка ещё велика, код дополнительно ослабляет \textit{faithfulness} и \textit{stability}. При целевом пороге
\begin{equation}
\tau_{\mathrm{main}} = 0.02
\end{equation}
и текущем значении $\mathcal{L}_{\mathrm{main}} > \tau_{\mathrm{main}}$ используется множитель
\begin{equation}
\xi = \frac{\tau_{\mathrm{main}}}{\mathcal{L}_{\mathrm{main}} + \varepsilon}.
\end{equation}
Тогда веса модифицируются как
\begin{equation}
(w_{\mathrm{cons}}, w_{\mathrm{spar}}, w_{\mathrm{faith}}, w_{\mathrm{stab}})
\leftarrow
\operatorname{Norm}
\bigl(
 w_{\mathrm{cons}}^{(0)},
 w_{\mathrm{spar}}^{(0)},
 \xi w_{\mathrm{faith}}^{(0)},
 \xi w_{\mathrm{stab}}^{(0)}
\bigr),
\end{equation}
где $\operatorname{Norm}$ означает повторную нормировку до суммы 1.

Если же \texttt{use\_adaptive\_weights=false}, то используются фиксированные веса из конфигурации:
\begin{equation}
(\gamma_{\mathrm{cons}}, \gamma_{\mathrm{spar}}, \gamma_{\mathrm{faith}}, \gamma_{\mathrm{stab}})
= (0.2, 0.7, 0.05, 0.05)
\end{equation}
с последующей нормировкой до единичной суммы. Однако для основной версии это запасной режим; в реальном финальном запуске активны именно адаптивные веса.

\section{Адаптивная балансировка SHAP и основной задачи}

Даже после агрегации компонент $\mathcal{L}_{\mathrm{SHAP}}$ дополнительно нормируется, чтобы интерпретируемость не подавляла основную задачу восстановления спектра.

\subsection{EMA основной ошибки}

В коде поддерживается экспоненциальное скользящее среднее основной ошибки:
\begin{equation}
M_t = \alpha M_{t-1} + (1-\alpha)\mathcal{L}_{\mathrm{main}}^{(t)},
\qquad
\alpha = 0.9.
\end{equation}

На первом шаге берётся $M_0 = \mathcal{L}_{\mathrm{main}}^{(0)}$.

\subsection{Динамическое целевое отношение SHAP/main}

Пусть $p_t \in [0,1]$ --- относительный прогресс по эпохам:
\begin{equation}
 p_t = \frac{t}{T},
\end{equation}
где $T$ --- общее число эпох второй стадии.

В коде целевое отношение SHAP/main выбирается как
\begin{equation}
\eta_t = \eta_0 (0.5 + 0.5 p_t),
\qquad
\eta_0 = 0.4.
\end{equation}

Таким образом, на ранних эпохах SHAP действует мягче, а к концу донастройки его допустимая доля в общей оптимизации увеличивается.

\subsection{Нормировка SHAP-штрафа}

Пусть
\begin{equation}
 r_t = \frac{\mathcal{L}_{\mathrm{SHAP}}^{(t)}}{M_t + \varepsilon}
\end{equation}
--- текущее отношение SHAP-штрафа к EMA основной ошибки.

Тогда в точной реализации используется масштабирование
\begin{equation}
 s_t =
 \begin{cases}
 \dfrac{\eta_t}{r_t}, & r_t > 2\eta_t,\\[6pt]
 \dfrac{\eta_t}{r_t}, & r_t < \eta_t/2,\\[6pt]
 1, & \eta_t/2 \le r_t \le 2\eta_t.
 \end{cases}
\end{equation}

То есть при слишком большом или слишком маленьком отношении SHAP к основной ошибке штраф нормируется к целевому диапазону.

\subsection{Замедление при отсутствии улучшения}

В тренере присутствует дополнительный множитель плавной сходимости. Если относительное улучшение основной ошибки становится меньше $0.1\%$, счётчик отсутствия улучшения увеличивается:
\begin{equation}
 q_t = q_{t-1} + 1.
\end{equation}
Тогда множитель замедления задаётся как
\begin{equation}
 c_t = \max\!\left(\frac{1}{1 + 0.1 q_t},\; c_{\min}\right),
 \qquad
 c_{\min} = 0.25.
\end{equation}

При заметном улучшении основной ошибки счётчик сбрасывается. Нормированный SHAP-штраф имеет вид
\begin{equation}
\widetilde{\mathcal{L}}_{\mathrm{SHAP}}^{(t)} = c_t s_t \mathcal{L}_{\mathrm{SHAP}}^{(t)}.
\end{equation}

\subsection{Расписание коэффициента gamma\_t}

Вес SHAP-регуляризации во внешнем функционале также изменяется в процессе донастройки. В коде используется кусочно-линейное расписание с параметрами
\begin{equation}
\gamma_{\mathrm{start}} = 0.02,
\qquad
\gamma_{\mathrm{end}} = 0.10,
\qquad
r_{\mathrm{warm}} = 0.5.
\end{equation}

Если $p_t = t/T$, то формула в текущей реализации имеет вид
\begin{equation}
\gamma_t =
\begin{cases}
\gamma_{\mathrm{start}} + (\gamma_{\mathrm{end}} - \gamma_{\mathrm{start}})\dfrac{p_t}{r_{\mathrm{warm}}}, & p_t < r_{\mathrm{warm}},\\[8pt]
\gamma_{\mathrm{end}}, & p_t \ge r_{\mathrm{warm}}.
\end{cases}
\end{equation}

На практике это означает мягкий линейный старт SHAP-регуляризации на первых $100r_{\mathrm{warm}}\%$ эпох и последующее удержание веса на уровне $\gamma_{\mathrm{end}}$.

\section{Тихоновская регуляризация гладкости спектра}

Главная версия проекта дополнительно использует тихоновскую регуляризацию по выходному спектру. Она вводится непосредственно на предсказаниях модели и штрафует резкие осцилляции между соседними энергетическими бинами.

Для одного предсказанного спектра
\begin{equation}
\hat{\mathbf{y}}_i = (\hat{y}_{i1},\dots,\hat{y}_{im})^\top
\end{equation}
определяются конечно-разностные операторы:
\begin{align}
(D_1 \hat{\mathbf{y}}_i)_k &= \hat{y}_{i,k+1} - \hat{y}_{i,k}, \qquad k=1,\dots,m-1,\\
(D_2 \hat{\mathbf{y}}_i)_k &= \hat{y}_{i,k+2} - 2\hat{y}_{i,k+1} + \hat{y}_{i,k}, \qquad k=1,\dots,m-2.
\end{align}

В проекте поддерживаются оба порядка, но основная рабочая версия использует второй порядок:
\begin{equation}
 p_{\mathrm{Tikh}} = 2.
\end{equation}

Соответствующий штраф по батчу записывается как
\begin{equation}
R_{\mathrm{Tikh}} = \frac{1}{B(m-2)} \sum_{i=1}^{B} \sum_{k=1}^{m-2}
\left(
\hat{y}_{i,k+2} - 2\hat{y}_{i,k+1} + \hat{y}_{i,k}
\right)^2.
\end{equation}

Именно эта формула соответствует реализации \texttt{mean(diffs**2)} в \texttt{shap\_trainer\_improved.py}. Второй порядок особенно важен для спектров, поскольку он штрафует не просто наклон, а локальную кривизну, то есть подавляет искусственные ``зубцы'' в восстановленном распределении.

\section{Итоговый функционал второй стадии}

После объединения всех членов вторая стадия минимизирует функционал
\begin{equation}
\mathcal{L}_{\mathrm{total}}^{(t)} =
\mathcal{L}_{\mathrm{main}}^{(t)}
+ \gamma_t \widetilde{\mathcal{L}}_{\mathrm{SHAP}}^{(t)}
+ \lambda_{\mathrm{Tikh}} R_{\mathrm{Tikh}}^{(t)}.
\end{equation}

В основной конфигурации
\begin{equation}
\lambda_{\mathrm{Tikh}} = 10^{-3},
\qquad
p_{\mathrm{Tikh}} = 2.
\end{equation}

Следовательно, ключевая версия метода представляет собой именно задачу
\begin{equation}
\theta^{\star} = \arg\min_{\theta}
\left[
\mathcal{L}_{\mathrm{main}}
+ \gamma_t \widetilde{\mathcal{L}}_{\mathrm{SHAP}}
+ 10^{-3} R_{\mathrm{Tikh}}
\right],
\end{equation}
где и интерпретируемость, и гладкость спектра являются мягкими регуляризаторами относительно основной задачи восстановления.

\section{Инференс}

После обучения модель используется как детерминированный оператор реконструкции:
\begin{equation}
\hat{\mathbf{y}} = f_{\theta^{\star}}(\mathbf{x}).
\end{equation}

Если данные были нормализованы по сумме, выполняется денормализация:
\begin{equation}
\hat{\mathbf{y}}_{\mathrm{denorm}} = S \hat{\mathbf{y}}'.
\end{equation}

В проекте дополнительно сохраняются массивы предсказаний, модель, важности признаков, история SHAP-компонент и графические артефакты для последующего анализа.

\section{Monte Carlo-оценка неопределённости}

Отдельный модуль проекта оценивает чувствительность прогноза к ошибкам во входных измерениях. Для заданного уровня относительной ошибки $\delta$ строятся возмущённые входы
\begin{equation}
\mathbf{x}^{(t)} = \mathbf{x} + \boldsymbol{\xi}^{(t)},
\qquad
\xi_j^{(t)} \sim \mathcal{N}\!\left(0, \sigma_j^2\right),
\qquad
\sigma_j = \frac{\delta}{100}|x_j|.
\end{equation}

В коде дополнительно выполняется отсечение отрицательных значений:
\begin{equation}
 x_j^{(t)} \leftarrow \max(x_j^{(t)}, 0).
\end{equation}

Если включена нормализация по сумме, то после внесения шума каждое возмущённое наблюдение снова нормируется:
\begin{equation}
\mathbf{x}^{(t)} \leftarrow \frac{\mathbf{x}^{(t)}}{\sum_j x_j^{(t)}}.
\end{equation}

После $T$ прогонов Monte Carlo получаем ансамбль прогнозов
\begin{equation}
\hat{\mathbf{y}}^{(t)} = f_{\theta^{\star}}(\mathbf{x}^{(t)}),
\qquad
t=1,\dots,T.
\end{equation}

По нему вычисляются:
\begin{align}
\bar{\mathbf{y}} &= \frac{1}{T} \sum_{t=1}^{T} \hat{\mathbf{y}}^{(t)},\\
\operatorname{Std}(\mathbf{y}) &= \sqrt{\frac{1}{T} \sum_{t=1}^{T} \left(\hat{\mathbf{y}}^{(t)} - \bar{\mathbf{y}}\right)^2},\\
\mathbf{y}^{(p)} &= \operatorname{Percentile}_p\left(\hat{\mathbf{y}}^{(1)},\dots,\hat{\mathbf{y}}^{(T)}\right), \quad p \in \{5,25,50,75,95\}.
\end{align}

Тем самым получаются среднее предсказание, стандартное отклонение, доверительные интервалы и распределение неопределённости по каждому энергетическому бину.

\section{Метрики качества}

Для финальной оценки на тестовом наборе используются следующие метрики.

\subsection{MSE, RMSE и MAE}

Для тестового множества из $N$ образцов:
\begin{equation}
\operatorname{MSE} = \frac{1}{N} \sum_{i=1}^{N} \frac{1}{m} \sum_{k=1}^{m} (\hat{y}_{ik} - y_{ik})^2,
\end{equation}
\begin{equation}
\operatorname{RMSE} = \sqrt{\operatorname{MSE}},
\end{equation}
\begin{equation}
\operatorname{MAE} = \frac{1}{N} \sum_{i=1}^{N} \frac{1}{m} \sum_{k=1}^{m} |\hat{y}_{ik} - y_{ik}|.
\end{equation}

\subsection{Взвешенный и средний \texorpdfstring{$R^2$}{R2} по бинам}

Поскольку задача многовыходная, в коде сначала считается отдельный коэффициент детерминации для каждого энергетического бина с ненулевой дисперсией:
\begin{equation}
R_k^2 = 1 - \frac{\sum_{i=1}^{N}(y_{ik} - \hat{y}_{ik})^2}{\sum_{i=1}^{N}(y_{ik} - \bar{y}_{k})^2},
\qquad
\bar{y}_k = \frac{1}{N} \sum_{i=1}^{N} y_{ik}.
\end{equation}

Далее вводятся веса, пропорциональные дисперсии бина:
\begin{equation}
\sigma_k^2 = \operatorname{Var}(y_{1k},\dots,y_{Nk}),
\qquad
\omega_k = \frac{\sigma_k^2}{\sum_{\ell} \sigma_{\ell}^2}.
\end{equation}

Тогда основная итоговая метрика в проекте --- variance-weighted $R^2$:
\begin{equation}
R^2_{\mathrm{weighted}} = \sum_{k} \omega_k R_k^2.
\end{equation}

Дополнительно сохраняется невзвешенное среднее по выходам:
\begin{equation}
R^2_{\mathrm{mean}} = \frac{1}{m'} \sum_{k=1}^{m'} R_k^2,
\end{equation}
где $m'$ --- число бинов с ненулевой дисперсией.

\subsection{Метрики по диапазонам энергий}

Для более детального анализа 60 бинов делятся на три диапазона:
\begin{equation}
[0,19], \qquad [20,39], \qquad [40,59].
\end{equation}
Для каждого диапазона отдельно вычисляются MSE, RMSE, MAE и $R^2$.

\section{Параметры основной Tikhonov-версии}

Финальная рабочая конфигурация метода задаётся следующими параметрами:
\begin{itemize}
    \item число правил ANFIS: $R=50$;
    \item функции принадлежности: Gaussian;
    \item этап 1 (PSO): 150 эпох, популяция 100;
    \item этап 2 (донастройка): 300 эпох, размер батча 32, $\mathrm{lr}=10^{-3}$;
    \item расписание SHAP-веса: $\gamma_{\mathrm{start}}=0.02$, $\gamma_{\mathrm{end}}=0.10$;
    \item warmup-доля для $\gamma_t$: $r_{\mathrm{warm}}=0.5$;
    \item adaptive gamma: включён;
    \item целевое отношение SHAP/main: $\eta_0=0.40$;
    \item минимальный множитель плавной сходимости: $c_{\min}=0.25$;
    \item adaptive component weights: включены;
    \item true-SHAP reference: включён;
    \item активные SHAP-компоненты: \textit{consistency}, \textit{sparsity}, \textit{faithfulness}, \textit{stability};
    \item резервные фиксированные веса компонент: $(0.2, 0.7, 0.05, 0.05)$ для \textit{consistency}, \textit{sparsity}, \textit{faithfulness}, \textit{stability};
    \item тихоновская регуляризация: $\lambda_{\mathrm{Tikh}}=2\cdot 10^{-3}$, порядок $p_{\mathrm{Tikh}}=2$.
\end{itemize}

\section{Практическая интерпретация метода}

С математической точки зрения итоговый метод сочетает три уровня априорной информации:
\begin{enumerate}
    \item \textbf{Аппроксимация данных:} член $\mathcal{L}_{\mathrm{main}}$ заставляет модель правильно восстанавливать спектр.
    \item \textbf{Интерпретируемость:} SHAP-компоненты добиваются того, чтобы карта важности детекторов была согласованной, разреженной, достоверной и устойчивой.
    \item \textbf{Физическая гладкость спектра:} тихоновский член подавляет нефизичные локальные осцилляции по энергии.
\end{enumerate}

По сути, оптимизация решает компромисс между точностью, объяснимостью и гладкостью, причём компромисс настраивается автоматически по ходу обучения через $\gamma_t$, EMA-нормировку и адаптивные веса компонент SHAP.

\section{Экспериментальные результаты для основной версии}

Для основного запуска \texttt{tikhonov\_stronger\_20260319} на финальном реальном тесте получены следующие метрики:
\begin{align}
    \operatorname{MSE} &= 0.01060296,\\
    \operatorname{RMSE} &= 0.10297068,\\
    \operatorname{MAE} &= 0.04890357,\\
    R^2_{\mathrm{weighted}} &= 0.83702629,\\
    R^2_{\mathrm{mean}} &= 0.55409377.
\end{align}

Метрики по диапазонам энергий:
\begin{itemize}
    \item бины $[0,19]$: $R^2_{\mathrm{weighted}} = 0.95486219$;
    \item бины $[20,39]$: $R^2_{\mathrm{weighted}} = 0.75988831$;
    \item бины $[40,59]$: $R^2_{\mathrm{weighted}} = 0.68916792$.
\end{itemize}

Нормализованная глобальная SHAP-важность в этом запуске образует распределение с единичной суммой.
Пять наиболее значимых детекторов имеют веса
\begin{equation}
Q3 = 0.269309,\qquad
Q5 = 0.206911,\qquad
Q6 = 0.137873,\qquad
Q4 = 0.095738,\qquad
Q1 = 0.078220.
\end{equation}
Это подтверждает, что модель не размывает вклад равномерно по всем 10 каналам, а выделяет устойчивое ядро информативных детекторов.

Средний вклад регуляризаторов в итоговый функционал для этого запуска составил
\begin{equation}
\mathbb{E}\,\bigl[\gamma_t \widetilde{\mathcal{L}}_{\mathrm{SHAP}}^{(t)}\bigr] = 5.54 \cdot 10^{-5},
\qquad
\mathbb{E}\,\bigl[\lambda_{\mathrm{Tikh}}\mathcal{L}_{\mathrm{Tikh}}^{(t)}\bigr] = 4.23 \cdot 10^{-4},
\end{equation}
а средняя доля регуляризации в полном функционале достигла
\begin{equation}
\mathbb{E}\,\left[
\frac{
\gamma_t \widetilde{\mathcal{L}}_{\mathrm{SHAP}}^{(t)}
\;+\;
\lambda_{\mathrm{Tikh}}\mathcal{L}_{\mathrm{Tikh}}^{(t)}
}{
\mathcal{L}_{\mathrm{total}}^{(t)}
}
\right]
= 0.0653.
\end{equation}

По сравнению с предыдущим более мягким режимом вклад SHAP во внешний функционал вырос примерно в $29.8$ раза, а вклад тихоновского члена --- примерно в $2.0$ раза, при этом итоговые метрики точности не ухудшились.

Дополнительный абляционный анализ на уменьшенной конфигурации показал:
\begin{itemize}
    \item удаление \textit{sparsity} даёт наибольшую просадку по $R^2_{\mathrm{weighted}}$, значит этот компонент является главным носителем полезной SHAP-структуры;
    \item удаление \textit{consistency} почти не ухудшает accuracy, но делает глобальную SHAP-карту более концентрированной и менее энтропийной;
    \item \textit{faithfulness} и \textit{stability} выступают как поддерживающие компоненты, влияющие слабее, но стабилизирующие итоговую структуру регуляризации.
\end{itemize}

При этом адаптивная схема весов выравнивает не сами коэффициенты, а \emph{эффективные взвешенные сигналы} компонент, поэтому даже при сильно различающихся raw weight значения
\textit{consistency}, \textit{sparsity}, \textit{faithfulness} и \textit{stability} дают сравнимый вклад в SHAP-часть функционала.

Для анализа чувствительности к входной погрешности был выполнен Monte Carlo-анализ. При уровне ошибки входных каналов $0.5\%$ получено:
\begin{equation}
\operatorname{mean\_std} = 3.8511 \cdot 10^{-5},
\qquad
\operatorname{max\_std} = 2.1553 \cdot 10^{-4}.
\end{equation}

При увеличении ошибки до $10\%$ соответствующие значения выросли до
\begin{equation}
\operatorname{mean\_std} = 7.5806 \cdot 10^{-4},
\qquad
\operatorname{max\_std} = 4.1703 \cdot 10^{-3}.
\end{equation}

Следовательно, главная \texttt{SHAP + Tikhonov} версия сохраняет приемлемую устойчивость к реалистичным ошибкам измерений и одновременно даёт физически более гладкие спектры по сравнению с вариантом без регуляризации гладкости.

\end{document}
